\chapter{Anosmia (Loss of Smell)}

\section*{Definition and Overview}
Anosmia is the complete loss of the sense of smell. Partial loss of smell is called hyposmia, and a distorted sense of smell, in which familiar odours are perceived differently and often unpleasantly, is called parosmia. Because the sense of smell contributes most of what is perceived as flavour, people with anosmia typically complain that food has lost its taste.

Loss of smell has many causes, ranging from a simple blocked nose to damage to the olfactory nerves or the brain. It became widely known during the COVID-19 pandemic as a common symptom of that infection. In most cases the condition is temporary and resolves as the underlying cause settles, but in a minority of people it persists or becomes permanent. Beyond the loss of enjoyment of food, anosmia has practical safety implications, because smell warns of smoke, gas, and spoiled food.

\section*{Epidemiology}
Loss of smell is very common and increases markedly with age. Roughly one in five older adults has some impairment of smell, and complete anosmia affects a smaller but significant proportion of the population. Upper respiratory tract infections are a leading cause of transient anosmia, and during the height of the COVID-19 pandemic very large numbers of people developed the symptom, most recovering within weeks.

Head injury, chronic nasal and sinus disease, nasal polyps, and ageing are other frequent causes. The sense of smell declines gradually with normal ageing, and men and women are affected at similar rates, although age is the strongest risk factor overall.

\section*{Aetiology and Pathophysiology}
Smell begins when odorant particles in the nose bind to specialised nerve cells in the olfactory epithelium, a small area high in the nasal cavity. Signals then travel along the olfactory nerve through the skull to the olfactory bulb and higher brain centres. Loss of smell can occur at any point in this pathway:

\begin{itemize}
    \item \textbf{Blockage or damage to the nasal lining:} chronic rhinosinusitis, nasal polyps, allergies, a deviated septum, or a cold -- the most common mechanism, usually reversible
    \item \textbf{Viral infection:} many viruses, including influenza and coronaviruses, can damage the olfactory nerves directly, even without nasal blockage
    \item \textbf{Head injury:} trauma can tear or crush the olfactory nerves where they cross the skull, and is a common cause in younger adults
    \item \textbf{Ageing:} the sense of smell naturally declines with age as olfactory nerve cells are lost
    \item \textbf{Neurological disease:} conditions such as Parkinson's disease, Alzheimer's disease, and other neurodegenerative disorders
    \item \textbf{Medications and toxins:} certain drugs and industrial chemicals can impair smell
    \item \textbf{Structural problems:} rarely, tumours or growths pressing on the olfactory pathway
    \item \textbf{Congenital causes:} a small number of people are born without a sense of smell, sometimes as part of a syndrome such as Kallmann syndrome
\end{itemize}

\section*{Clinical Features}
The main feature is reduced or absent smell, but associated symptoms vary with the cause:

\begin{itemize}
    \item Inability to detect odours, including warning smells such as smoke or gas
    \item Loss of food flavour, since much of what is perceived as taste is actually smell
    \item Reduced interest in eating and possible unintentional weight loss
    \item With nasal causes: a blocked or runny nose, facial pressure, or a fluctuating sense of smell
    \item With viral anosmia: sudden loss of smell, sometimes with no other cold symptoms
    \item After head injury: loss of smell accompanied by the effects of the injury itself
    \item In some people, parosmia, in which everyday smells such as coffee or onions are perceived as foul or distorted
\end{itemize}

\section*{Differential Diagnosis}
The key clinical distinction is between causes affecting the nasal passages (conductive) and those affecting the nerves or brain (sensorineural):

\begin{itemize}
    \item \textbf{Chronic rhinosinusitis, with or without polyps:} the most common treatable cause
    \item \textbf{Post-viral anosmia:} including after COVID-19, influenza, and the common cold
    \item \textbf{Head injury:} with or without other neurological symptoms
    \item \textbf{Neurodegenerative disease:} early loss of smell may be the first sign of Parkinson's or Alzheimer's disease
    \item \textbf{Nasal tumours:} rare, usually associated with one-sided blockage or bleeding
    \item \textbf{Medication-related loss:} caused by certain drugs, including some antibiotics and blood pressure medicines
    \item \textbf{Congenital anosmia:} present since birth, including in Kallmann syndrome, where it is accompanied by delayed puberty
\end{itemize}

\section*{When to Seek Medical Care}
See a doctor if loss of smell persists for more than a few weeks, follows a head injury, or affects only one side of the nose. One-sided or progressive symptoms should always be assessed.

\textbf{Seek urgent care if:}
\begin{itemize}
    \item Loss of smell is accompanied by other neurological symptoms such as weakness, numbness, or difficulty speaking
    \item There is a severe or persistent headache, especially after a head injury
    \item There is bleeding from the nose or a blocked feeling on one side only
    \item A child fails to develop a sense of smell or has delayed puberty along with anosmia
\end{itemize}

\textbf{Call 000 (or local emergency services) for:}
\begin{itemize}
    \item Sudden neurological symptoms such as facial drooping, arm weakness, or difficulty speaking, which may indicate a stroke
    \item Loss of consciousness or a major head injury
\end{itemize}

\section*{Diagnosis and Investigations}
A doctor will take a history and examine the nose, and may arrange further tests:

\begin{itemize}
    \item \textbf{History:} when the loss began, whether it is one-sided or both-sided, and any recent colds, injury, or medication changes
    \item \textbf{Examination of the nose:} with an endoscope (a thin, flexible camera) to inspect the nasal passages and olfactory region
    \item \textbf{Smell testing:} validated scratch-and-sniff tests to measure the ability to detect and identify odours
    \item \textbf{Imaging:} a CT scan of the sinuses for chronic sinus disease, or an MRI of the brain and olfactory region when a nerve or brain cause is suspected
    \item \textbf{Blood tests:} occasionally used to look for nutritional, hormonal, or inflammatory causes
    \item \textbf{Referral:} to an ear, nose, and throat (ENT) specialist or a neurologist for persistent or unexplained cases
\end{itemize}

\section*{Management}
Treatment depends on the cause:

\begin{itemize}
    \item \textbf{Nasal causes:} steroid sprays or drops, saline nasal rinses, and treatment of underlying allergy or sinus infection
    \item \textbf{Nasal polyps or structural blockage:} surgery such as polypectomy or functional endoscopic sinus surgery can restore airflow and smell
    \item \textbf{Post-viral anosmia:} most cases recover spontaneously over weeks to months, so the main approach is supportive
    \item \textbf{Smell training:} deliberately sniffing strong, familiar scents such as rose, lemon, eucalyptus, and cloves twice daily for several months can encourage recovery of the olfactory nerves
    \item \textbf{Stopping offending medications:} under medical advice, when a drug is the cause
    \item \textbf{Safety measures:} installing smoke alarms, checking gas appliances, and dating food to reduce the risks of not smelling hazards
\end{itemize}

\section*{Complications}
\begin{itemize}
    \item \textbf{Impaired safety:} inability to smell smoke, gas leaks, or spoiled food
    \item \textbf{Reduced enjoyment of food:} poor appetite and unintentional weight loss
    \item \textbf{Missed warning signs:} inability to detect one's own body odours or the smell of burning
    \item \textbf{Psychological distress:} loss of smell can be isolating and is associated with low mood and depression
    \item \textbf{Underlying disease:} anosmia may be an early clue to a progressive neurological condition
\end{itemize}

\section*{Prognosis and Outcomes}
The outlook depends on the cause. Anosmia from a blocked nose or sinus infection usually resolves when the underlying problem is treated. Post-viral anosmia, including after COVID-19, resolves in the majority of people, although recovery may take months and some people are left with persistent or distorted smell. Loss of smell from ageing is usually gradual and permanent. Where the nerves are damaged by head injury, recovery is variable and may be incomplete. Smell training appears to improve the chances of recovery in many cases, so it is worth persisting with it.

\section*{Prevention and Self-Care}
\begin{itemize}
    \item Treat colds and sinus infections promptly, and manage allergies to keep the nose clear
    \item Avoid prolonged exposure to toxic fumes, and use ventilation or masks when working with chemicals
    \item Wear protective headgear during sport and use seatbelts to reduce the risk of head injury
    \item If anosmia is present, install smoke detectors, label foods with dates, and have gas appliances regularly checked
    \item Try smell training with familiar scents twice daily for several months
    \item Report any new neurological symptoms or one-sided symptoms to a doctor
    \item Use flavour, texture, and colour to make food more appealing if the sense of taste is affected
\end{itemize}

\section*{Sources and Further Reading}
The following sources provide reliable further information on anosmia:

\begin{itemize}
    \item NHS. \textit{Information on loss of smell and its causes and treatments.} https://www.nhs.uk/conditions/loss-of-smell/
    \item Mayo Clinic. \textit{Overview of smell disorders, including causes and treatment.} https://www.mayoclinic.org/diseases-conditions/smell-symptoms/basics/definition/sym-20050904
    \item MedlinePlus. \textit{Anosmia and smell disorders: overview and treatment.} https://medlineplus.gov/smelldisorders.html
    \item MSD Manual. \textit{Smell and taste disorders in the consumer health section.} https://www.msdmanuals.com/home/ear-nose-and-throat-disorders/taste-and-smell-disorders
    \item Centers for Disease Control and Prevention. \textit{Information on COVID-19 symptoms, including loss of smell.} https://www.cdc.gov/coronavirus/2019-ncov/symptoms-testing/symptoms.html
\end{itemize}
